\chapter{功与机械能}\label{ch:g12:work-and-energy}

把一颗弹珠投入沙拉碗中：它先下沉，越过碗底，
再爬回起始高度，稍作停顿，然后折返——每一次
折返都可以由一条曲线（碗的剖面）和一条水平线
（弹珠的能量）预测出来，无需求解运动方程。
去年在直线上计算了力的功（\cref{ch:g11:work-of-force}），
并比较了动能与势能
（\cref{ch:g11:mechanical-energy}）；今年要把这套核算推广到
各种情形：曲线路径、弹簧，以及一目了然的势能地形。

\section{沿任意路径的功}

\begin{definition}[沿路径的功]\label{def:g12:work-and-energy:path}
设力 $\vect F$ 作用于沿任一路径从 $A$ 运动到
$B$ 的物体上：把路径分成足够短的位移 $\vect{\dd\ell}$，
每段近似为直线，且力 $\vect F$ 在其上几乎不变。每一小段
做的元功 $\scal{F}{\dd\ell} = F\,\dd\ell\cos\theta$
（\cref{ch:g11:work-of-force}），
\emph{沿路径的功}\index{功（力的）!沿路径的功}即为
各元功之和 $W_{AB}(\vect F) = \sum \scal{F}{\dd\ell}$——在力恒定的直线路径上，它
便退化为去年的
$F \times AB \times \cos\theta$。
\end{definition}

\begin{omfigure}
\begin{tikzpicture}[font=\small]
  \draw[omExo!60, thick] plot [smooth, tension=0.8] coordinates {(0,0.3) (1.2,1.15) (2.6,1.3) (4.0,0.7) (5.3,1.05)};
  \draw[omMeth, thick] (0,0.3) -- (1.2,1.15) -- (2.6,1.3) -- (4.0,0.7) -- (5.3,1.05);
  \foreach \p in {(0,0.3),(1.2,1.15),(2.6,1.3),(4.0,0.7),(5.3,1.05)} \fill \p circle (1.3pt);
  \node[below] at (0,0.3) {$A$}; \node[below right] at (5.3,1.05) {$B$};
  \draw[-Stealth, omThm, very thick] (2.6,1.3) -- (3.44,0.94) node[below left=-2pt] {$\vect{\dd\ell}$};
  \draw[-Stealth, omDef, very thick] (2.6,1.3) -- ++(22:1.5) node[above] {$\vect F$};
  \draw (2.6,1.3) ++(-23:0.75) arc[start angle=-23, end angle=22, radius=0.75];
  \node at ($(2.6,1.3)+(0:1.0)$) {$\theta$};
\end{tikzpicture}

{\small 曲线路径是由许多短小的折线段组成的：每一段
$\vect{\dd\ell}$ 做功 $F\,\dd\ell\cos\theta$，总功即各段之和。}
\end{omfigure}

\begin{proposition}[重力所做的功，任意路径]\label{prop:g12:work-and-energy:weight}
沿\emph{任一}路径，从 $A$（高度 $z_A$）到 $B$（高度
$z_B$），质量为 $m$ 的物体所受重力做的功为
$W_{AB}(\vect P) = mg\,(z_A - z_B)$：仅取决于下降的高度，与路径无关。
\end{proposition}

\begin{proof}
每一小段贡献 $mg \times (\text{微小下降量})$（去年
的直线段计算）；各段的下降量逐项相消，
得 $z_A - z_B$。去年承认的曲线情形，现在得到证明。
\end{proof}

\begin{definition}[保守力与非保守力]\label{def:g12:work-and-energy:conservative}
称力为\emph{保守力}\index{保守力}，若其在两点之间所做的功
沿每条路径都相同（绕任一闭合路径一周时功为零），
否则称为\emph{非保守力}\index{非保守力}
。重力与弹簧力（见下一节）是保守力。
滑动摩擦力则不然：它与运动方向相反，做功 $W = -fL$，
即在长为 $L$ 的路径上——每米收取一次“过路费”，且永不退还。
\end{definition}

\begin{example}[通往小屋的两条路线]\label{ex:g12:work-and-energy:trails}
一名质量为 \qty{75}{kg} 的徒步者，沿一条长 \qty{800}{m} 的陡峭小径
或 \qty{2.4}{km} 长的盘山路，都可上升 \qty{300}{m}。重力
在两条路线上做功 $-75 \times 9.81 \times 300 \approx \qty{-2.2e5}{J}$；
一个 \qty{30}{N} 的摩擦力，在一条路线上做功 \qty{-2.4e4}{J}，
在另一条上为 \qty{-7.2e4}{J}。重力在下山时把功还了回来；摩擦力的
代价都化作热量散失，无论往返皆然。
\end{example}

\section{势能}

\begin{definition}[势能]\label{def:g12:work-and-energy:ep}
对每一个保守力，都引入一个
\emph{势能}\index{势能!由保守力做功}
$E_p$，一个仅取决于位置的数，定义式为
$W_{AB} = E_p(A) - E_p(B)$：所做之功即为所消耗的势能。
在某一选定的参考点处取 $E_p = 0$，即可处处确定其值；
参考点的不同选取只会使所有数值平移同一常量，
该常量在任何差值中都相互抵消。对重力而言，这就得到去年的 $E_p = mgz$，
现在则适用于任意路径。摩擦力不能拥有 $E_p$：与路径有关的
代价不是两端点数值之差——只有保守力
能够可取回地储存，这恰是“保守”二字所指的含义。
\end{definition}

\begin{proposition}[弹性势能]\label{prop:g12:work-and-energy:elastic}
一根劲度系数为 $k$（\unit{N/m}）的弹簧，被拉伸或压缩 $x$
后，以力 $F = kx$ 向回拉
（\cref{ch:g12:oscillators-and-time}）；弹簧回到静止时所做的功
——即其
\emph{弹性势能}\index{势能!弹性}——为
$E_p = \tfrac12\, k x^2$。
\end{proposition}

\begin{proof}
回复力沿运动方向作用，因此 $\sum F\,\dd\ell$ 即为
\emph{$F$--$x$ 图像下的面积}：一个底为 $x$、
高为 $kx$ 的三角形，面积为 $\tfrac12 x \times kx$。同一伸长量处的力
去程与回程相同：功只由两端点决定，弹簧力为
保守力，$E_p$ 有良好定义。
\end{proof}

\begin{omfigure}
\begin{tikzpicture}
\begin{axis}[omaxis, width=8.4cm, height=4.8cm, xmin=0, xmax=0.095,
    ymin=0, ymax=21,
    xlabel={伸长量 $x$ (\unit{m})}, ylabel={力 $F$ (\unit{N})}]
  \addplot[draw=none, fill=omDef!18, domain=0:0.06, samples=2] {200*x} \closedcycle;
  \addplot[omThm, thick, domain=0:0.09, samples=2] {200*x};
  \draw[dashed, omExo] (axis cs:0,12) -- (axis cs:0.06,12) -- (axis cs:0.06,0);
  \node[font=\small, anchor=west] at (axis cs:0.061,6.5) {$kx$};
  \node[font=\small] at (axis cs:0.044,3.2) {$\tfrac12 kx^2$};
  \node[font=\small, anchor=south, omThm] at (axis cs:0.082,17) {$F = kx$};
\end{axis}
\end{tikzpicture}

{\small 弹簧的力随伸长量增大；储存的能量
即直线下方的面积：一个三角形，$\tfrac12 \times x \times kx$。}
\end{omfigure}

\section{两个能量定理}

\begin{theorem}[动能定理]\label{thm:g12:work-and-energy:ketheorem}\index{动能定理}
对质量为 $m$ 的物体从 $A$ 到 $B$ 的任意运动——曲线的、
回环的、三维的——
\[
\Delta E_k = E_k(B) - E_k(A) = \sum W_{AB}(\vect F),
\]
其中求和遍及作用于物体的所有力。
\end{theorem}

\begin{proof}
由于 $v^2 = \scal{v}{v}$，由本学年数学中的乘积法则
可得 $v^2$ 的导数为 $2\,\scal{a}{v}$，因此
$\dd E_k/\dd t = m\,\scal{a}{v} = \scal{F_{\text{总}}}{v}$，由
牛顿第二定律（见\cref{ch:g12:newtons-laws}）。在很短的 $\dd t$
时间内物体移动 $\vect{\dd\ell} = \vect v\,\dd t$，$E_k$ 获得
（见\cref{def:g12:work-and-energy:path}）的微功；把各步相加，
即各功相加——上一学年承认的结论，现已完整证得。
\end{proof}

\begin{proposition}[瞬时功率]\label{prop:g12:work-and-energy:power}
任一时刻，力 $\vect F$ 作用于速度为 $\vect v$ 的物体上，
其功率为 $P = \scal{F}{v} = Fv\cos\theta$，且
$\dd E_k/\dd t$ 等于所有力的功率之和：上一学年的
匀速运动公式，对任意运动的每一瞬时均严格成立。
\end{proposition}

\begin{proof}
由前一证明可读出：每个力所做的功
$\scal{F}{\dd\ell} = \scal{F}{v}\,\dd t$ 计入 $E_k$，发生于 $\dd t$ 内。
\end{proof}

\begin{theorem}[机械能定理]\label{thm:g12:work-and-energy:emtheorem}\index{机械能定理}
设 $E_m = E_k + E_p$，其中 $E_p$ 为
所有起作用的保守力的势能之和。沿任一路径，
\[
\Delta E_m = W_{AB}(\text{非保守力}),
\]
且当非保守力不做功时，$E_m$ 守恒。
\end{theorem}

\begin{proof}
将动能定理拆开：$\Delta E_k = W_{\text{保守}} +
W_{\text{非保守}} = -\Delta E_p + W_{\text{非保守}}$；将 $\Delta E_p$ 移项。
\end{proof}

\begin{example}[单摆，这次认真算一算]\label{ex:g12:work-and-energy:pendulum}
一个用长为 $L = \qty{2.5}{m}$ 的细线悬挂的摆球，
由静止释放，离竖直方向 $\qty{45}{\degree}$。张力
与运动方向垂直，不做功；重力为保守力：
$E_m$ 在\emph{曲线圆弧上}守恒，这一结论上一学年只能
验证。摆球上升的高度 $h = L(1 - \cos\qty{45}{\degree}) =
\qty{0.73}{m}$，于是在最低点处 $v = \sqrt{2gh} \approx
\qty{3.8}{m/s}$。
\end{example}

\begin{method}[任意运动的能量分析]\label{met:g12:work-and-energy:audit}
\begin{enumerate}
  \item 选取研究系统，并列出其所受的力。
  \item 分类：与运动垂直 $\to$ 不做功；保守力
        $\to$ 以 $E_p$ 计入 $E_m$；其余的 $\to$ $W_{\text{非保守}}$
        （滑动摩擦力：$-fL$，$L$ 为路径长度）。
  \item 在已知点与所求点之间写出 $\Delta E_m = W_{\text{非保守}}$，
        求解。能量可回答“多快，多
        高”；只有时间和加速度才需要牛顿定律。
\end{enumerate}
\end{method}

\section{能量图}

\begin{definition}[能量图，转折点]\label{def:g12:work-and-energy:diagram}
对沿单一坐标轴 $x$ 运动且 $E_m$ 守恒的运动，
\emph{能量图}\index{能量图}画出曲线 $E_p(x)$
与水平线 $E_m$。由于 $E_k = E_m - E_p \geq 0$，
运动被限制在曲线位于直线下方的区域内；两者之间的间隙
即为 $E_k$，可直接读出。在
\emph{转折点}\index{转折点}处，即曲线与直线相交处，
$v = 0$，运动方向反转。
\end{definition}

\begin{omfigure}
\begin{tikzpicture}
\begin{axis}[omaxis, width=9.0cm, height=5.4cm,
    xlabel={$x$}, ylabel={能量}, xtick=\empty, ytick=\empty,
    xmin=-2.25, xmax=2.25, ymin=0, ymax=2.45]
  \addplot[omDef, thick, domain=-2.05:2.05, samples=80] {0.25*x^4 - x^2 + 1.5};
  \addplot[omThm, dashed, thick, domain=-2.25:2.25, samples=2] {1.0};
  \addplot[omProp, dashed, domain=-2.25:2.25, samples=2] {1.7};
  \node[font=\small, omThm, anchor=west] at (axis cs:-2.2,0.86) {$E_m$};
  \node[font=\small, omProp, anchor=west] at (axis cs:-2.2,1.85) {$E_m'$};
  \node[font=\small, omDef, anchor=west] at (axis cs:1.55,2.1) {$E_p(x)$};
  \fill (axis cs:0.765,1.0) circle (1.6pt) node[font=\small, below=1pt] {$x_1$};
  \fill (axis cs:1.848,1.0) circle (1.6pt) node[font=\small, below right=-1pt] {$x_2$};
  \fill[omMeth] (axis cs:1.414,0.5) circle (2.2pt) node[font=\small, below] {稳定的};
  \fill[omMeth] (axis cs:0,1.5) circle (2.2pt) node[font=\small, above] {不稳定的};
\end{axis}
\end{tikzpicture}

{\small 能量图，无需求解即可读出：
能量为 $E_m$ 时，物体在 $x_1$ 与 $x_2$ 之间振动，间隙最大处速率最快；
能量为 $E_m'$ 时，它缓慢爬过山丘并逃逸。}
\end{omfigure}

\begin{proposition}[由势能曲线得力；平衡]\label{prop:g12:work-and-energy:equilibria}
沿 $x$ 方向的保守力等于 $E_p$ 斜率的相反数：
$F = -\,\dd E_p/\dd x$——在图上指向\emph{下坡方向}，
且在切线水平处为零，此处即为一个平衡。
$E_p$ 的极小值处为\emph{稳定平衡}\index{稳定平衡}——偏离后的
物体被推回，如同谷底的小球；极大值处则为
\emph{不稳定平衡}\index{不稳定平衡}——物体被推离，
如同山顶的小球。
\end{proposition}

\begin{proof}
在微小位移 $\dd x$ 上，力所做的功为 $F\,\dd x = E_p(x) -
E_p(x + \dd x) = -\dd E_p$；除以 $\dd x$。谷底两侧的下坡方向
都指向谷内，山顶附近则指向山外。
\end{proof}

\begin{remark}[振动与逃逸]\label{rem:g12:work-and-energy:escape}
\cref{ch:g12:oscillators-and-time}中的微小振动，都位于
$E_p$ 曲线的谷底。远离地球时，引力减弱，
$E_p(r)$ 曲线越来越缓慢地攀升，趋向一个有限的上限（
即渐弱的力图线下的面积是有限的）。若火箭的 $E_m$
达到该上限，便不会遇到转折点，它
\emph{逃逸}\index{逃逸速度}——从地面发射时，速率约为 \qty{11.2}{km/s}
（见\cref{exo:g12:work-and-energy:13}）；低于此值，它仍被束缚。
该上限本身在本科一年级分册中算出。
\end{remark}

\begin{example}[竖直圆环]\label{ex:g12:work-and-energy:loop}
一颗弹珠从高度 $H$ 处释放，滑入一个
半径为 $R$ 的竖直圆环。圆周运动的动力学（见\cref{ch:g12:newtons-laws}）要求
最高点处 $v_{\text{最高}}^2 \geq gR$——这是去年过山车问题借用的数据，如今轮到我们了。
由机械能守恒，从静止的 $H$ 处到达最高点（高度 $2R$）
可得 $v_{\text{最高}}^2 = 2g(H - 2R)$，于是
$H \geq \tfrac52 R$：比圆环最高点高出半个半径，与
质量无关——在有摩擦的真实世界中，所需高度还需略高一些。
\end{example}

\begin{omfigure}
\begin{tikzpicture}[font=\small, scale=0.9]
  \draw[thick] (0,3.0) .. controls (0.9,1.1) and (1.6,0) .. (3.1,0) -- (6.9,0);
  \draw[thick] (5.5,1.2) circle (1.2);
  \fill[omDef] (0,3.0) circle (2.2pt);
  \draw[dashed, omExo] (4.0,2.4) -- (7.6,2.4) node[right] {$2R$};
  \draw[dashed, omThm] (-0.4,3.0) -- (7.6,3.0) node[right] {$H_{\min} = \tfrac52 R$};
  \draw[Stealth-Stealth, omThm] (7.35,2.4) -- (7.35,3.0);
  \node[right, omThm] at (7.35,2.7) {$\tfrac{R}{2}$};
  \draw[-Stealth, omMeth, thick] (5.5,2.4) -- (4.5,2.4) node[above right=0pt and 2pt] {$v_{\text{最高}}$};
  \fill (5.5,2.4) circle (1.5pt);
\end{tikzpicture}

{\small 竖直圆环：释放点高出圆环最高点半个半径，
故剩余的 $E_k$ 使弹珠仍压在轨道上（$v_{\text{最高}}^2 = gR$）。}
\end{omfigure}

\section{练习}

\begin{exercise}[$\star$]\label{exo:g12:work-and-energy:1}
一根劲度系数为 $k = \qty{200}{N/m}$ 的弹簧，被拉伸
\qty{5.0}{cm}。求其储存的能量；若拉伸量加倍，
则能量变为多少？\emph{第二段} \qty{5.0}{cm} 的拉伸需做多少功？
\end{exercise}

\begin{exercise}[$\star$]\label{exo:g12:work-and-energy:2}
一架 \qty{0.90}{kg} 的无人机沿蜿蜒航线，从露台
（$z = \qty{12}{m}$）飞往阳台（$z = \qty{30}{m}$）。重力做的功为多少？
沿直线爬升呢？往返一周时呢？
\end{exercise}

\begin{exercise}[$\star$]\label{exo:g12:work-and-energy:3}
判断下列各力是保守力、非保守力，还是不做功？请从路径出发说明理由：
(a) 重力；(b) 弹簧力；(c) 滑动摩擦力；(d) 滑动箱子所受的
法向力；(e) 空气阻力。
\end{exercise}

\begin{exercise}[$\star$]\label{exo:g12:work-and-energy:4}
一个摆长 \qty{2.0}{m} 的单摆，从与竖直方向
成 $\qty{60}{\degree}$ 角处静止释放，求其在最低点的速率；
为何细线的张力从不进入能量守恒的计算？
\end{exercise}

\begin{exercise}[$\star$]\label{exo:g12:work-and-energy:5}
一名骑车人在平路上以 \qty{9.0}{m/s} 的速率骑行，输出
\qty{250}{W} 的功率：她所对抗的总阻力为多大？若以相同功率爬坡，
速率为 \qty{5.0}{m/s}，她所克服的力为多大？
\end{exercise}

\begin{exercise}[$\star\star$]\label{exo:g12:work-and-energy:6}
玩具发射器的弹簧（$k = \qty{450}{N/m}$）被压缩
\qty{8.0}{cm}，紧靠一枚 \qty{20}{g} 的飞镖。求储存的能量、
发射速率，以及若竖直向上发射（不计空气阻力）所能达到的高度。
\end{exercise}

\begin{exercise}[$\star\star$]\label{exo:g12:work-and-energy:7}
拉伸橡皮筋时，一名学生测得 $F = 0$，$3.0$，$7.0$，
$12.0$，\qty{18.0}{N}，对应 $x = 0$，$2.0$，$4.0$，$6.0$，\qty{8.0}{cm}。
用梯形面积估算储存的能量，并与
$\tfrac12 kx^2$ 比较（取末端的 $k = F/x$）。两者为何不同？
\end{exercise}

\begin{exercise}[$\star\star$]\label{exo:g12:work-and-energy:8}
一名质量 \qty{55}{kg} 的雪橇者从静止出发，沿
长 \qty{25}{m} 的曲线滑道下降 \qty{8.0}{m}，到达时速率为 \qty{9.0}{m/s}。
机械能损失了多少？平均摩擦力为多大？为何滑道的具体形状无关紧要？
\end{exercise}

\begin{exercise}[$\star\star$]\label{exo:g12:work-and-energy:9}
一条弹珠轨道含有半径为 \qty{20}{cm} 的圆环。求
最小释放高度（由静止释放，无摩擦）。若从 \qty{60}{cm} 处释放，
求圆环最高点的速率，以及比值 $v_{\text{最高}}^2/(gR)$？为何
实际轨道必须从更高处释放？
\end{exercise}

\begin{exercise}[$\star\star$]\label{exo:g12:work-and-energy:10}
一颗 \qty{100}{g} 的珠子沿某一坐标轴无摩擦滑动，
已知 $E_p(x) = 8.0\,x^2$（以焦耳计，$x$ 以米为单位），$E_m = \qty{2.0}{J}$。
求转折点？最大速率为多少？这是何种运动？
\end{exercise}

\begin{exercise}[$\star\star$]\label{exo:g12:work-and-energy:11}
一辆 \qty{1200}{kg} 的汽车以 \qty{2.0}{m/s} 的速率，被
劲度系数为 $k = \qty{6.0e5}{N/m}$ 的缓冲弹簧制停。求最大压缩量。
速率为 \qty{4.0}{m/s} 时，为何压缩量并非四倍，尽管能量变为四倍？
\end{exercise}

\begin{exercise}[$\star\star\star$]\label{exo:g12:work-and-energy:12}
一辆 \qty{0.50}{kg} 的小车在轨道上滚动，该轨道的 $E_p(x)$ 曲线有一个
\qty{6.0}{J} 的峰顶（$x = 0$）和一个 \qty{2.0}{J} 的谷底
（$x = \qty{1.5}{m}$）。(a) 小车从左侧驶来，
机械能 $E_m = \qty{5.0}{J}$，能否越过峰顶？若不能，将如何运动？(b) 若
$E_m = \qty{7.0}{J}$ 时，峰顶与谷底的速率各为多少？(c) 平衡位置有哪些？
\end{exercise}

\begin{exercise}[$\star\star\star$]\label{exo:g12:work-and-energy:13}
在本阶段不加证明地承认，将质量为 $m$ 的物体从地面移至任意远处
需做功 $mgR_E$，$R_E = \qty{6.37e6}{m}$（即弱化引力图像下的
有限面积），求地球的\emph{逃逸速率}。
为何探测器与小石子的结果相同？若发射速率略低于它，又会怎样？
\end{exercise}

\begin{exercise}[$\star\star\star$]\label{exo:g12:work-and-energy:14}
一个 \qty{250}{g} 的滑块在弹簧（$k = \qty{40}{N/m}$）上振动，
振幅为 \qty{6.0}{cm}（见\cref{ch:g12:oscillators-and-time}）。
总能量是多少？最大速率是多少？$E_k = E_p$ 的位置在何处？
\end{exercise}

\begin{exercise}[$\star\star\star$]\label{exo:g12:work-and-energy:15}
一名撑杆跳运动员以 \qty{9.5}{m/s} 的速率冲刺。若撑杆将其全部动能转化为
势能，试估算其质心升高的高度，以及所能越过的横杆高度
（质心起始高度为 \qty{1.0}{m}）。世界纪录约为
\qty{6.3}{m}：该估算计入了什么，又遗漏了什么？
\end{exercise}

\section{问题：跳台滑雪}

\begin{problem}[{周末问题——跳台滑雪：经测速枪核验的助滑道、
自由落体式的飞行、陡坡着陆的缓和几何，
以及摩擦力买单后设计者不得不建造的跳台}]
\label{pb:g12:work-and-energy:1}
一名奥运会跳台滑雪运动员（$m = \qty{70}{kg}$，含滑雪板）
从高出起跳台 $H = \qty{45}{m}$ 的起滑门处静止出发，
滑过长为 $L = \qty{90}{m}$ 的弯曲助滑道，并沿水平方向离开起跳台；
测速枪在此读数为 $v_0 = \qty{26}{m/s}$。着陆点
比起跳台低 $h = \qty{40}{m}$，位于倾角为 $\qty{35}{\degree}$ 的斜坡上。第 II
、III 部分忽略空气阻力。

\medskip\noindent\textbf{第一部分——助滑道。}
\begin{enumerate}
  \item 预测无摩擦时的起跳速率；助滑道的
        曲线形状为何无关紧要？
  \item 以起跳台为参考面，计算起滑门处与起跳时的 $E_m$。
  \item 计算 $\Delta E_m$；哪条定理指明了其中的元凶，该定理出自何人？
  \item 推算滑行全程的平均阻力（雪面与空气阻力之和）。
  \item $E_m$ 损失了多大比例？为何要给滑雪板打蜡、屈膝下蹲？
\end{enumerate}

\medskip\noindent\textbf{第二部分——飞行。}
\begin{enumerate}[resume]
  \item 飞行中仅受重力作用，这属于哪种运动（见
        \cref{ch:g12:projectile-motion}）？着陆时
        速度分量各为多少？
  \item 着陆速率需算两次：由分量求得，并由 $E_m$
        守恒求得。以 \unit{km/h} 表示是多少？
  \item 计算飞行时间及水平飞行距离。
  \item 实际跳台滑雪运动员借助空气如同机翼般滑翔，飞行距离远得多：
        空气能否\emph{提高}着陆速率？请运用
        机械能定理加以论证。
\end{enumerate}

\medskip\noindent\textbf{第三部分——着陆。}
\begin{enumerate}[resume]
  \item 计算着陆速度在水平线以下的夹角。
  \item 将着陆速度分解为平行于斜坡与垂直于斜坡的
        两个分量。
  \item 双腿只吸收垂直于斜坡的那部分：计算该分量
        对应的动能，并与总动能相比较。
  \item 若为平地，动能将全部被吸收：各次着陆
        等效于竖直下落多高（$h_{\text{等效}} = v_\perp^2/2g$，
        分别为 $v^2/2g$）？
  \item 用一句话说明：着陆坡为何既陡峭又呈曲线？
\end{enumerate}

\medskip\noindent\textbf{第四部分——设计者的跳台。}
较小的跳台只需 $v_0' = \qty{24}{m/s}$ 的起跳速率；其
直线助滑道的倾角为 $\qty{35}{\degree}$，因而起滑门高度 $H'$
即助滑道长 $H'/\sin\qty{35}{\degree}$；阻力同样为 \qty{80}{N}。
\begin{enumerate}[resume]
  \item 计算无摩擦情形下设计者应建造的起滑门高度。
  \item 对 $H'$ 写出从起滑门到起跳台的机械能定理。
  \item 求解 $H'$。
  \item 摩擦力使跳台的高度提高了百分之几？
  \item 此时质量不再约去：对一名质量为
        \qty{60}{kg} 的跳台滑雪运动员重新计算 $H'$。谁需要更高的跳台？为什么？
  \item 点睛之笔：宣布的跳台高度，以及其中摩擦力所占的份额。
\end{enumerate}
\end{problem}
